%% Separately authored additions after exact Intuitionistic Logic restoration.
% Necessity: the exact inherited chapter does not contain the registered IPL
% system card, two countermodel figures, Kolmogorov/Kuroda variants and
% comparison table, or the later algebraic/topological/first-order/intermediate
% material, so those separately attributed reader aids would otherwise vanish.
% Definition: this editorial derivative descends from the six marked historical
% ranges but updates captions, consumers, system standing, and later material;
% altered ranges are derivatives rather than verbatim slices.
% Standing: it is editorial reader material, not inherited source or a silent
% proof-status promotion.
% Nonredundancy: the exact wrappers provide the source chapter and its proofs;
% these slices retain only registered apparatus and later topics, while the
% historical combined body and annex remain byte-pinned and unreachable.
% Use: parts/p08.tex inputs this file exactly once after SRC-13180-13200; the
% route, slice, label, consumer, and publication tests fail closed on drift.
\noindent\textit{Reader note.} The following system card, figures, translation
comparisons, and later semantics are separately attributed additions to
the exact inherited chapter above; they neither replace nor alter its claims.

\paragraph{Preorder semantics for the algebraic constructions.}
The inherited presentation uses partial orders. The following separately
stated extension permits nonempty preorders, using the same persistent
valuation and forcing clauses
\cite{Kripke1965,TroelstraSchwichtenberg2000}.
For these definitions, let \(\mathsf{Prop}\) be any specified set of
propositional variables, with the connective grammar of
Section~\ref{sec:intuitionistic-logic-01}.

\begin{hbdefinition}{def:p08-ipl-kripke-structure}{Intuitionistic Kripke frame and structure on a preorder}
An intuitionistic Kripke frame is a nonempty preordered set
\((W,\leq)\): the relation is reflexive and transitive.
An intuitionistic Kripke structure is
\(\mathcal M=(W,\leq,V)\), where
\(V:\mathsf{Prop}\to\mathcal P(W)\) is persistent:
for every \(p\in\mathsf{Prop}\),
\[
              w\in V(p),\quad w\leq v
                     \quad\Longrightarrow\quad v\in V(p).
\]
Partial orders are the antisymmetric special case.
The preorder is the persistence relation; an interpretation as growth
of information does not add a semantic clause.
\end{hbdefinition}

\begin{hbdefinition}{def:p08-ipl-forcing}{Intuitionistic forcing on a preorder}
For a structure \(\mathcal M=(W,\leq,V)\) and \(w\in W\),
define forcing recursively.
\begin{itemize}
\item \(\mathcal M,w\forces p\) exactly when \(w\in V(p)\);
      no world forces \(\bottom\).
\item \(\mathcal M,w\forces A\land B\) exactly when it forces both
      \(A\) and \(B\).
\item \(\mathcal M,w\forces A\lor B\) exactly when it forces at least
      one of \(A,B\) at that same world.
\item \(\mathcal M,w\forces A\to B\) exactly when, for every
      \(v\geq w\), if \(\mathcal M,v\forces A\), then
      \(\mathcal M,v\forces B\).
\end{itemize}
Here \(\neg A\) abbreviates \(A\to\bottom\).
The conditional in the last clause is metalinguistic; the
object-language connective occurs inside the formula being evaluated.
\end{hbdefinition}

\begin{hbdefinition}{def:p08-ipl-consequence}{Local intuitionistic consequence on nonempty preorders}
For a set \(\Gamma\) of formulas and a formula \(A\), write
\(\Gamma\entails_{\mathsf{IPL}}A\) exactly when, for every
structure of Definition~\ref{def:p08-ipl-kripke-structure}
and every world \(w\), forcing each member of \(\Gamma\) at \(w\)
implies forcing \(A\) at \(w\).
This is semantic consequence at a common world.
The finite-support derivability relation of
Section~\ref{sec:intuitionistic-logic-01} remains a distinct
syntactic relation.
\end{hbdefinition}

% BEGIN HISTORICAL DERIVATIVE 90-128
\begin{systementry}{entry:p08-ipl}{Intuitionistic Propositional Logic}
  \entryaliases{The inherited v1.1 presentation uses \(\mathsf{IPL}\).
    \(\mathsf{IPC}\) is a literature alias for this same propositional logic,
    not a second system.}
  \entrylanguage{Formulas over
    \(\bottom,\land,\lor,\to\), with
    \(\neg A\defeq A\to\bottom\), as fixed in
    Section~\ref{sec:intuitionistic-logic-01}.}
  \entrysemantics{Persistent valuations on the partially ordered Kripke frames fixed in
    Section~\ref{sec:intuitionistic-logic-01}, with forcing and consequence
    as defined there.}
  \entrycalculus{sheet:lj}
  \entryadequacy{proved=annex:p8.ipl.source}{For every formula set \(\Gamma\)
    and formula \(A\),
    \(\Gamma\derives_{\mathsf{IPL}}A\) exactly when
    \(\Gamma\entails_{\mathsf{IPL}}A\).  If finite \(\Gamma\) does not
    derive \(A\), a finite rooted Kripke structure has a world forcing every
    member of \(\Gamma\) but not \(A\).}
  \entrydecidability{proved=annex:p8.ipl.decision-card}{Derivability in the
    propositional fragment is decidable; finite filtration gives an effective
    finite search bound over the subformulas of the input sequent.}
  \entrycharacteristic{Excluded middle fails in
    Figure~\ref{fig:p08-ipl-excluded-middle}, as does double-negation
    elimination; Peirce's law fails in
    Figure~\ref{fig:p08-ipl-peirce}.  The full disjunction property and its
    cut-free last-rule component are recorded, without duplication, in
    Theorem~\ref{thm:p07-ipc-disjunction-property} and
    Lemma~\ref{thm:p07-lj-disjunction-last-rule}.}
  \entrytranslations{For this card, \texttt{sheet:lj} means exactly its
    propositional subcalculus: omit \(\forall L/R\) and \(\exists L/R\),
    retain the displayed propositional, structural, Cut, and Mix rows, and keep
    the single-succedent restriction.  See Theorems~\ref{thm:glivenko}
    and~\ref{thm:negative-translation}, and the four-map comparison in
    Table~\ref{tab:p08-negative-translations}.}
  \entryalgebra{Heyting-algebra semantics and algebraic completeness are
    stated in Theorem~\ref{thm:p08-heyting-completeness}.}
  \entrysources{Kripke1965,TroelstraVanDalen1988,TroelstraSchwichtenberg2000}
  \index{intuitionistic propositional logic}
  \index{IPC|see{intuitionistic propositional logic}}
  \index{IPL|see{intuitionistic propositional logic}}
\end{systementry}
% END HISTORICAL DERIVATIVE 90-128

% BEGIN HISTORICAL DERIVATIVE 149-165
\begin{figure}[!htbp]
\centering
\begin{tikzpicture}[>=Latex,node distance=3.3cm]
  \node[draw,rounded corners,align=center,inner sep=6pt] (w0)
    {\(w_0\)\\\(\mathcal M,w_0\nforces p\)\\
     \(\mathcal M,w_0\nforces\neg p\)\\
     \(\mathcal M,w_0\nforces p\lor\neg p\)};
  \node[draw,rounded corners,align=center,inner sep=6pt,above=of w0] (w1)
    {\(w_1\)\\\(\mathcal M,w_1\forces p\)\\
     \(\mathcal M,w_1\nforces\neg p\)};
  \draw[->] (w0) -- node[right]{\(w_0<w_1\)} (w1);
\end{tikzpicture}
\caption{The two-world persistent model refuting excluded middle at \(w_0\).
The annotations state forcing at each world rather than
  treating eventual forcing as disjunction at the root.}
\label{fig:p08-ipl-excluded-middle}
\end{figure}
% END HISTORICAL DERIVATIVE 149-165

% BEGIN HISTORICAL DERIVATIVE 175-191
\begin{figure}[!htbp]
\centering
\begin{tikzpicture}[>=Latex,node distance=3.5cm]
  \node[draw,rounded corners,align=center,inner sep=6pt] (w0)
    {\(w_0\)\\\(\mathcal M,w_0\forces((p\to q)\to p)\)\\
     \(\mathcal M,w_0\nforces p\)\\
     \(\mathcal M,w_0\nforces(((p\to q)\to p)\to p)\)};
  \node[draw,rounded corners,align=center,inner sep=6pt,above=of w0] (w1)
    {\(w_1\)\\\(\mathcal M,w_1\forces p\)\\
     \(\mathcal M,w_1\nforces q\)\\
     \(\mathcal M,w_1\nforces p\to q\)};
  \draw[->] (w0) -- node[right]{\(w_0<w_1\)} (w1);
\end{tikzpicture}
\caption{A two-world intuitionistic countermodel to Peirce's law
\(((p\to q)\to p)\to p\).}
\label{fig:p08-ipl-peirce}
\end{figure}
% END HISTORICAL DERIVATIVE 175-191

% BEGIN HISTORICAL DERIVATIVE 208-275
\begin{hbdefinition}{def:p08-negative-translation-maps}{Glivenko, Kolmogorov, and Kuroda maps}
Fix a first-order formula grammar with primitive \(\bottom\).  In this
definition \(P\) ranges only over genuine atomic formulas, so it never denotes
\(\bottom\), and \(\neg A\) remains the abbreviation
\(A\to\bottom\).  The following expressions name metalinguistic maps on
formulas, not new object-language connectives.

The Glivenko map is the nonrecursive propositional map
\(G(A)\defeq\neg\neg A\).  The Kolmogorov translation is the recursion
\[
 \begin{aligned}
 P^{\mathrm{Ko}}&\defeq\neg\neg P,
 &\bottom^{\mathrm{Ko}}&\defeq\neg\neg\bottom,\\
 (A\land B)^{\mathrm{Ko}}&\defeq
   \neg\neg\bigl(A^{\mathrm{Ko}}\land B^{\mathrm{Ko}}\bigr),
 &(A\lor B)^{\mathrm{Ko}}&\defeq
   \neg\neg\bigl(A^{\mathrm{Ko}}\lor B^{\mathrm{Ko}}\bigr),\\
 (A\to B)^{\mathrm{Ko}}&\defeq
   \neg\neg\bigl(A^{\mathrm{Ko}}\to B^{\mathrm{Ko}}\bigr),
 &(\forall x\,A)^{\mathrm{Ko}}&\defeq
   \neg\neg\forall x\,A^{\mathrm{Ko}},\\
 (\exists x\,A)^{\mathrm{Ko}}&\defeq
   \neg\neg\exists x\,A^{\mathrm{Ko}}.
 \end{aligned}
\]
For the Kuroda translation, first define the inner map by
\[
 \begin{aligned}
 P_{\mathrm{Ku}}&\defeq P,
 &\bottom_{\mathrm{Ku}}&\defeq\bottom,\\
 (A\land B)_{\mathrm{Ku}}&\defeq A_{\mathrm{Ku}}\land B_{\mathrm{Ku}},
 &(A\lor B)_{\mathrm{Ku}}&\defeq A_{\mathrm{Ku}}\lor B_{\mathrm{Ku}},\\
 (A\to B)_{\mathrm{Ku}}&\defeq A_{\mathrm{Ku}}\to B_{\mathrm{Ku}},
 &(\exists x\,A)_{\mathrm{Ku}}&\defeq\exists x\,A_{\mathrm{Ku}},\\
 (\forall x\,A)_{\mathrm{Ku}}&\defeq
   \forall x\,\neg\neg A_{\mathrm{Ku}},
 \end{aligned}
\]
and then put \(A^{\mathrm{Ku}}\defeq\neg\neg A_{\mathrm{Ku}}\).
For either recursive translation, a superscript on a context means that every
member is translated.

Separate names are necessary because the maps place double negations at
different construction stages: Glivenko wraps only a whole propositional
formula, Kolmogorov wraps every construction result, and Kuroda modifies the
universal clause before applying one outer double negation.  Those differences
are used in the two cited embedding theorems and in
Table~\ref{tab:p08-negative-translations}.
\index{negative translation!Glivenko}
\index{negative translation!Kolmogorov}
\index{negative translation!Kuroda}
\end{hbdefinition}

\begin{hbtheorem}{thm:p08-kolmogorov-embedding}{cited=Kolmogorov1925ExcludedMiddle}{Kolmogorov negative-translation theorem}
For every finite first-order formula context \(\Gamma\) and formula \(A\),
\(\Gamma\derives_{\mathsf{LK}}A\) exactly when
\(\Gamma^{\mathrm{Ko}}\derives_{\mathsf{LJ}}A^{\mathrm{Ko}}\).
In particular, the theoremhood form translates the conclusion and has no
untranslated premise context.
\end{hbtheorem}

\begin{hbtheorem}{thm:p08-kuroda-embedding}{cited=Kuroda1951}{Kuroda negative-translation theorem}
For every finite first-order formula context \(\Gamma\) and formula \(A\),
\(\Gamma\derives_{\mathsf{LK}}A\) exactly when
\(\Gamma^{\mathrm{Ku}}\derives_{\mathsf{LJ}}A^{\mathrm{Ku}}\).
The target is intuitionistic logic; no corresponding minimal-logic claim is
made here.
\end{hbtheorem}
% END HISTORICAL DERIVATIVE 208-275

% BEGIN HISTORICAL DERIVATIVE 313-379
\begin{landscape}
\begingroup\scriptsize\setlength{\tabcolsep}{3pt}
\begin{longtable}{@{}p{.75in}p{3.60in}p{2.85in}@{}}
\caption{Negative translations into intuitionistic logic.}
\label{tab:p08-negative-translations}\\
\toprule
Map & Exact clauses & Preservation statement and evidence \\
\midrule
\endfirsthead
\caption[]{Negative translations into intuitionistic logic (continued)}\\
\toprule
Map & Exact clauses & Preservation statement and evidence \\
\midrule
\endhead
\midrule\multicolumn{3}{r}{\emph{continued on next page}}\\
\endfoot
\endlastfoot
Glivenko \(G\) &
Not structural: \(G(A)\defeq\neg\neg A\), so in particular
\(G(\bottom)=\neg\neg\bottom\). &
Propositional rules only: the exact finite-context equivalence is
Theorem~\ref{thm:glivenko} --- \texttt{cited=Glivenko1929}.  No
quantifier extension is asserted. \\
\addlinespace
Kolmogorov \(\mathrm{Ko}\) &
\(P^{\mathrm{Ko}}\defeq\neg\neg P\) and
\(\bottom^{\mathrm{Ko}}\defeq\neg\neg\bottom\).  Each binary result is
double-negated:
\((A\land B)^{\mathrm{Ko}}\defeq
\neg\neg(A^{\mathrm{Ko}}\land B^{\mathrm{Ko}})\),
\((A\lor B)^{\mathrm{Ko}}\defeq
\neg\neg(A^{\mathrm{Ko}}\lor B^{\mathrm{Ko}})\), and
\((A\to B)^{\mathrm{Ko}}\defeq
\neg\neg(A^{\mathrm{Ko}}\to B^{\mathrm{Ko}})\).  Likewise,
\((\forall x\,A)^{\mathrm{Ko}}\defeq
\neg\neg\forall x\,A^{\mathrm{Ko}}\) and
\((\exists x\,A)^{\mathrm{Ko}}\defeq
\neg\neg\exists x\,A^{\mathrm{Ko}}\). &
Classical \(\mathsf{LK}\) derivability is equivalent to
\(\mathsf{LJ}\) derivability after translating the conclusion and every
premise: Theorem~\ref{thm:p08-kolmogorov-embedding} ---
\texttt{cited=Kolmogorov1925ExcludedMiddle}. \\
\addlinespace
G\"odel--Gentzen \(\mathrm N\) &
\(\bottom^{\mathrm N}\defeq\bottom\), with priority over
\(P^{\mathrm N}\defeq\neg\neg P\).  The map commutes with
\(\land,\to,\forall\), while
\((A\lor B)^{\mathrm N}\defeq
\neg(\neg A^{\mathrm N}\land\neg B^{\mathrm N})\) and
\((\exists x\,A)^{\mathrm N}\defeq
\neg\forall x\,\neg A^{\mathrm N}\). &
The exact \(\mathsf{LK}\)-to-\(\mathsf{LJ}\) sequent equivalence, with
memberwise context translation, is Theorem~\ref{thm:negative-translation} ---
\texttt{cited=Godel1933}; see also \cite{Gentzen1936}. \\
\addlinespace
Kuroda \(\mathrm{Ku}\) &
The inner map has \(P_{\mathrm{Ku}}\defeq P\) and
\(\bottom_{\mathrm{Ku}}\defeq\bottom\), commutes with
\(\land,\lor,\to,\exists\), and satisfies
\((\forall x\,A)_{\mathrm{Ku}}\defeq
\forall x\,\neg\neg A_{\mathrm{Ku}}\).  The full map is
\(A^{\mathrm{Ku}}\defeq\neg\neg A_{\mathrm{Ku}}\), so
\(\bottom^{\mathrm{Ku}}=\neg\neg\bottom\). &
Classical \(\mathsf{LK}\) derivability is equivalent to intuitionistic
\(\mathsf{LJ}\) derivability after translating the conclusion and every
premise: Theorem~\ref{thm:p08-kuroda-embedding} ---
\texttt{cited=Kuroda1951}.  No minimal-logic target is claimed. \\
\bottomrule
\end{longtable}
\endgroup
\end{landscape}
% END HISTORICAL DERIVATIVE 313-379

% BEGIN HISTORICAL DERIVATIVE 396-681
\section{Algebraic and topological semantics}

\begin{hbdefinition}{def:p08-heyting-algebra}{Heyting algebra and formula valuation}
A \emph{Heyting algebra} is a bounded distributive lattice

\[
  \mathcal H=(H,\wedge,\vee,0,1,\Rightarrow)
\]

whose final binary operation is characterized, for all \(a,b,c\in H\), by

\[
  a\wedge b\leq c
  \quad\text{exactly when}\quad
  a\leq b\Rightarrow c.
\]

In this definition the source symbols \(\wedge,\vee,0,1,\Rightarrow\) are
typed as Heyting-algebra operations and bounds; in particular,
\(\Rightarrow\) is not the sequent arrow and not the object-language
conditional.  A formula valuation is a map
\(\nu:\mathsf{Prop}\to H\), extended by
\(\nu(\bottom)=0\), by the lattice meet and join for
\(\land,\lor\), and by
\(\nu(A\to B)=\nu(A)\Rightarrow\nu(B)\).
Algebraic negation is
\(\neg a\defeq a\Rightarrow0\).

This structure is needed to distinguish algebraic validity from Kripke
forcing.  Boolean algebras are the special Heyting algebras in which every
\(a\vee\neg a=1\); general Heyting
algebras need not satisfy that equation.  The distinction is used by the
IPL/CPL comparison and the intermediate-logics table below.
\index{Heyting algebra}
\index{semantics!algebraic!intuitionistic}
\end{hbdefinition}

\begin{hbtheorem}{thm:p08-heyting-completeness}{cited=RasiowaSikorski1963}{Heyting-algebra completeness for IPL}
For every propositional formula \(A\),
\(\derives_{\mathsf{IPL}}A\) exactly when every valuation in every Heyting
algebra assigns \(A\) the top value \(1\).
\end{hbtheorem}

The theorem is a theoremhood statement.  No arbitrary-premise algebraic
consequence relation is implied without separately fixing filter semantics.
The full algebraic proof and its finite-premise meet-inequality form
are Theorem~\ref{thm:p10-heyting-completeness}.
Prime-filter representation and strong local Kripke completeness,
with their choice assumptions explicit, are developed in
Section~\ref{sec:p10-heyting-completeness}.

\begin{hbdefinition}{def:p08-ipl-topological-semantics}{Open-set semantics for IPL}
Let \(X\) be a topological space.  An \emph{open-set valuation} is a map
\(\nu:\mathsf{Prop}\to\mathcal P(X)\) whose values are open subsets of
\(X\), extended to formulas by

\[
\begin{aligned}
 \nu(\bottom)&\defeq\varnothing,
 &\nu(A\land B)&\defeq\nu(A)\cap\nu(B),\\
 \nu(A\lor B)&\defeq\nu(A)\cup\nu(B),
 &\nu(A\to B)&\defeq
   \operatorname{int}_{X}
   \bigl((X\setminus\nu(A))\cup\nu(B)\bigr).
\end{aligned}
\]

Thus every formula denotes an open set.  The interior operation belongs to the
semantic algebra of opens; it is not a new object-language connective.  A
formula is valid under this semantics exactly when its value is \(X\) under
every open-set valuation.
\index{semantics!topological!intuitionistic}
\index{open-set valuation}
\end{hbdefinition}

\begin{hbtheorem}{thm:p08-topological-completeness}{cited=McKinseyTarski1946}{Topological completeness for IPL}
The theorems of IPL are exactly the formulas valid under all open-set
valuations on all topological spaces.
\end{hbtheorem}

\section{Constructive readings and boundaries}

Brouwer's 1908 paper supplies the historical critique of unrestricted logical
principles, especially excluded middle; it is not cited as the source of the
later six-clause BHK package \cite{Brouwer1908}.

\begin{hbdefinition}{def:p08-bhk-reading}{Brouwer--Heyting--Kolmogorov explanatory reading}
The label \emph{BHK reading} is a metalinguistic abbreviation for the
following explanatory clauses, not the name of another derivability or
satisfaction relation.  A construction of \(A\land B\) is a pair of
constructions, one for each conjunct; a construction of \(A\lor B\) is a
tagged choice of a disjunct together with its construction; and a construction
of \(A\to B\) is a method transforming each construction of \(A\) into one of
\(B\).  There is no construction of \(\bottom\).  A construction of
\(\forall x\,A(x)\) is a method producing a construction of \(A(a)\) for every
admissible \(a\); a construction of \(\exists x\,A(x)\) supplies an admissible
witness together with a construction of its instance
\cite{Heyting1956,TroelstraVanDalen1988}.

The name is needed only to retrieve these six clauses as one explanatory
package.  It differs from both the formally recursive forcing relation in
Definition~\ref{sec:intuitionistic-logic-01} and the number-realizability relation in
Definition~\ref{def:p07-number-realizability}; the latter assigns natural
numbers and program indices according to an exact clause table.  This
distinction is used when reading the cited realizability theorem below.
\index{Brouwer--Heyting--Kolmogorov interpretation}
\index{BHK|see{Brouwer--Heyting--Kolmogorov interpretation}}
\end{hbdefinition}

The formal result relevant to realizability is
Theorem~\ref{thm:p07-number-realizability-soundness}: from an HA derivation of
a closed arithmetic sentence it effectively produces a realizing number, with
status \texttt{cited=Kleene1945}.  The BHK clauses alone are neither that
theorem nor a soundness or completeness theorem for IPL.

\begin{proofboundary}{def:p08-intuitionistic-mathematics-boundary}{From IPL to formal intuitionistic mathematics}
Choice sequences, bar-induction principles, and continuity principles live in
formal theories of intuitionistic analysis, not in the propositional calculus
of \hyperref[entry:p08-ipl]{the IPL card}.  Their exact schemas depend on a
particular source theory, sequence language, and admissible bar or functional
class.  Until those data are fixed, a philosophical description cannot carry
a derivability, validity, or independence status.  The countermodels above
are Brouwerian-style counterexamples to particular classical formulas; they do
not by themselves axiomatize intuitionistic analysis.
\index{proof boundary!intuitionistic mathematics}
\end{proofboundary}

\begin{hbdefinition}{def:p08-choice-sequence-principles}{Choice-sequence, bar-induction, and continuity principles}
Work in the two-sorted language of numbers and sequences
\(\alpha:\mathbb N\to\mathbb N\), with finite initial segments
\(\bar\alpha n\). For a decidable predicate \(B\) on finite sequences, write
\(\operatorname{Bar}(B)\) for
\(\forall\alpha\,\exists n\,B(\bar\alpha n)\). The decidable bar-induction
scheme is
\[
 \operatorname{Bar}(B)\land
 \forall s(B(s)\to C(s))\land
 \forall s((\forall m\,C(s^\frown m))\to C(s))
 \ \to\ C(\langle\rangle).
\]
The continuity principle for a total functional
\(F:\mathbb N^{\mathbb N}\to\mathbb N\) is
\[
 \forall\alpha\,\exists n\,\forall\beta
 (\bar\alpha n=\bar\beta n\to F(\alpha)=F(\beta)).
\]
Choice sequences are the sequence-sort objects; no decidability of their
future values is assumed. The package name is necessary because BHK meaning,
number realizability, bar induction, and continuity are distinct formal
commitments. It stands as displayed axiom schemes in intuitionistic analysis,
differs from HA by its sequence sort and principles, and is used only under
these exact hypotheses \cite{TroelstraVanDalen1988}.
\index{choice sequence}\index{bar induction}\index{continuity principle!intuitionistic analysis}
\end{hbdefinition}

Markov's principle is not duplicated here. Its exact arithmetic scheme,
number-realizability validation, and nonderivability boundary for unrestricted
HA are Definition~\ref{def:p07-markov-principle} and
Theorem~\ref{thm:p07-markov-realizability-boundary}; those results do not imply
the choice-sequence or continuity schemes above.

\begin{hbtheorem}{thm:p08-normal-existence-extraction}{cited=TroelstraSchwichtenberg2000}{Existence extraction from a closed normal derivation}
Fix the pure intuitionistic first-order natural-deduction calculus on
\hyperref[sheet:nj]{\texttt{sheet:nj}}, with no equality rules, nonlogical
axioms, or free parameters.  If it has a closed normal derivation, with no open
assumptions, of \(\exists x\,A(x)\), then its final inference supplies a closed
term \(t\) and a derivation of \(A(t)\).
\end{hbtheorem}

This is a calculus-bounded normal-form statement.  The source normalization
theorem and its exact proof-sketch boundary are in
Section~\ref{sec:normalization-curry-howard-01}.  Nothing here
promotes the distinct HA numerical-existence statement or adds equality and
arithmetic induction to \texttt{sheet:nj}.

\section{First-order intuitionistic logic}

\begin{hbdefinition}{def:p08-iqc-expanding-domain}{Expanding-domain IQC model}
Here \emph{IQC} names the pure first-order logic of
\hyperref[sheet:lj]{\texttt{sheet:lj}} under the semantics now defined; the
name is needed to distinguish this setting from propositional IPL, but it does
not select a second rule sheet.  Fix an equality-free signature \(\Lang\).
An \emph{expanding-domain IQC model} consists of a nonempty preorder
\((W,\leq)\), a nonempty domain \(D_w\) for every \(w\in W\), and
interpretations at the worlds, subject to the following persistence data.
If \(w\leq u\), then \(D_w\subseteq D_u\); every constant denotes the same
element at later worlds; every function interpretation at \(u\) extends its
interpretation at \(w\); and every predicate extension at \(w\) is included in
its extension at \(u\).  These conditions make term values stable when an
assignment \(s:\mathsf{Var}\to D_w\) is viewed at a later world.

Atomic forcing uses the predicate interpretation at \(w\), and the
propositional clauses are those of
Definition~\ref{sec:intuitionistic-logic-01}.  The quantifier clauses are

\begin{itemize}
\item \(\mathcal M,w,s\forces\exists x\,A\) exactly when some
  \(a\in D_w\) has
  \(\mathcal M,w,s[x\mapsto a]\forces A\);
\item \(\mathcal M,w,s\forces\forall x\,A\) exactly when, for every
  \(u\geq w\) and every \(a\in D_u\),
  \(\mathcal M,u,s[x\mapsto a]\forces A\).
\end{itemize}

The expanding-domain hypothesis is load-bearing for persistence.  These
clauses define a satisfaction relation for IQC models, not modal
quantification and not a new object-language quantifier.
\index{intuitionistic first-order logic}
\index{Kripke model!expanding-domain intuitionistic}
\index{IQC|see{intuitionistic first-order logic}}
\end{hbdefinition}

\begin{hbtheorem}{thm:p08-iqc-completeness}{cited=Kripke1965}{IQC Kripke completeness}
For every closed equality-free first-order formula \(A\),
\(\derives_{\mathsf{LJ}}A\) exactly when \(A\) is forced at every world under
every assignment in every expanding-domain IQC model of
Definition~\ref{def:p08-iqc-expanding-domain}.
\end{hbtheorem}

\begin{hbdefinition}{def:p08-iqc-constant-domain-principle}{Constant-domain principle}
For a formula \(A\), a formula \(B(x)\), and a variable \(x\) not free in
\(A\), the \emph{constant-domain principle} is

\[
  \forall x\,(A\lor B(x))\to
  \bigl(A\lor\forall x\,B(x)\bigr).
\]

Its standing is an additional first-order intuitionistic schema.  It is not a
modal Barcan formula: the displayed quantifier is governed by the IQC clauses
above, and no modal operator occurs.
\index{constant-domain principle}
\end{hbdefinition}

\begin{hbtheorem}{thm:p08-iqc-constant-domain-boundary}{cited=TroelstraVanDalen1988}{Constant-domain validity boundary}
Every instance of the schema in
Definition~\ref{def:p08-iqc-constant-domain-principle} is valid when all
worlds have one common domain.  The schema can fail on expanding-domain IQC
models and is not a theorem of IQC.
\end{hbtheorem}

\section{Intermediate logics}

\begin{hbdefinition}{def:p08-intermediate-logic}{Intermediate propositional logic}
An \emph{intermediate propositional logic} is a set \(L\) of propositional
formulas which contains every IPL theorem, is contained in classical
propositional logic, and is closed under uniform substitution and modus
ponens.  This theorem-set notion is needed to compare strict extensions of
IPL without identifying them with arbitrary Kripke frame classes.  It differs
from an individual calculus presentation; the two instances below are used in
Table~\ref{tab:p08-intermediate-logics}.
\index{intermediate logic}
\index{superintuitionistic logic|see{intermediate logic}}
\end{hbdefinition}

\begin{hbdefinition}{def:p08-intermediate-lattice}{The complete lattice of intermediate logics}
Order intermediate logics by theorem-set inclusion:
\[
 L\leq M\quad\Longleftrightarrow\quad L\subseteq M.
\]
For any family \(\{L_i:i\in I\}\), its meet is
\(\bigcap_{i\in I}L_i\).  Its join is the least intermediate logic containing
\(\bigcup_{i\in I}L_i\), equivalently the closure of that union under uniform
substitution and modus ponens.  With the empty-family conventions, the meet of
no logics is \(\mathsf{CPC}\) and the join of no logics is \(\mathsf{IPC}\).
Thus all intermediate logics form a complete lattice with bottom
\(\mathsf{IPC}\) and top \(\mathsf{CPC}\).

The lattice language is necessary because most intermediate logics are not
linearly ordered by inclusion.  It stands as a metalinguistic order on theorem
sets, not as a new object-language connective, and differs from the small
strict chain displayed below.  That chain is only a local subdiagram of the
full superintuitionistic lattice
\cite{ChagrovZakharyaschev1997}.
\index{intermediate logic!lattice}
\index{superintuitionistic logic!lattice}
\end{hbdefinition}

\begin{hbtheorem}{thm:p08-kc-profile}{cited=Jankov1968WeakExcludedMiddle}{KC profile}
The intermediate logic \(\mathsf{KC}\) is IPL extended by the schema
\(\neg A\lor\neg\neg A\).  It has the finite-model property; its standard
finite rooted-frame presentation uses a unique greatest world.  It does not
have the theorem-level disjunction property.
\end{hbtheorem}

\begin{hbtheorem}{thm:p08-lc-profile}{cited=Dummett1959}{G\"odel--Dummett logic profile}
The intermediate logic \(\mathsf{LC}\) is IPL extended by the schema
\((A\to B)\lor(B\to A)\).  It is characterized by linearly ordered
intuitionistic frames, has the finite-model property, and does not have the
theorem-level disjunction property.
\end{hbtheorem}

In each system the added schema is itself a theorem-shaped disjunction while
neither displayed disjunct is a theorem of that system; hence the ordinary
theorem-level disjunction property would fail.  This does not contradict the
IPL property cited in Theorem~\ref{thm:p07-ipc-disjunction-property}.

The two logics are ordered, not incomparable.  Substitute \(B=\neg A\) in
LC prelinearity.  IPL proves
\[
 (A\to\neg A)\to\neg A,
 \qquad
 (\neg A\to A)\to\neg\neg A;
\]
disjunction elimination therefore derives
\(\neg A\lor\neg\neg A\).  Hence every KC theorem is an LC theorem.
The inclusion is strict: take the rooted frame
\[
 r<a<t,\qquad r<b<t,
\]
where \(a,b\) are incomparable and \(t\) is the unique greatest world, and put
\(V(p)=\{a,t\}\), \(V(q)=\{b,t\}\).  This is a KC frame, but at \(r\)
\(p\to q\) fails at the extension \(a\) and \(q\to p\) fails at \(b\).
Thus LC prelinearity fails there.  The outer inclusions are strict as well:
a root with two terminal successors, one forcing \(p\) and the other forcing
\(\neg p\), refutes \(\neg p\lor\neg\neg p\); and a two-world chain whose top
alone forces \(p\) refutes \(p\lor\neg p\) at its root.  Hence
\[
 \mathsf{IPC}\subsetneq\mathsf{KC}\subsetneq
 \mathsf{LC}\subsetneq\mathsf{CPC}.
\]

\begingroup\scriptsize
\begin{longtable}{@{}p{.12\textwidth}p{.26\textwidth}p{.28\textwidth}p{.13\textwidth}p{.07\textwidth}@{}}
\caption{Intermediate logics KC and LC.}
\label{tab:p08-intermediate-logics}\\
\toprule
System & Added schema over IPL & Frame presentation & Disjunction\newline property & FMP \\
\midrule
\endfirsthead
\caption[]{Intermediate logics KC and LC (continued)}\\
\toprule
System & Added schema over IPL & Frame presentation & Disjunction\newline property & FMP \\
\midrule
\endhead
\midrule\multicolumn{5}{r}{\emph{continued on next page}}\\
\endfoot
\endlastfoot
\(\mathsf{KC}\) & \(\neg A\lor\neg\neg A\) & Finite rooted frames with a
unique greatest world; cited profile in
Theorem~\ref{thm:p08-kc-profile} & No & Yes \\
\(\mathsf{LC}\) & \((A\to B)\lor(B\to A)\) & Linearly ordered
intuitionistic frames; cited profile in
Theorem~\ref{thm:p08-lc-profile} & No & Yes \\
\bottomrule
\end{longtable}
\endgroup

The table records the two systems separately; the derivation and countermodel
above establish the strict inclusion \(\mathsf{KC}\subsetneq\mathsf{LC}\).

\begin{figure}[htbp]
\centering
\begin{tikzpicture}[>=Stealth,node distance=16mm]
  \node[draw,rounded corners] (ipc) {\(\mathsf{IPC}\)};
  \node[draw,rounded corners,above=of ipc] (kc) {\(\mathsf{KC}\)};
  \node[draw,rounded corners,above=of kc] (lc) {\(\mathsf{LC}\)};
  \node[draw,rounded corners,above=of lc] (cpc) {\(\mathsf{CPC}\)};
  \draw[->] (ipc) -- (kc);
  \draw[->] (kc) -- (lc);
  \draw[->] (lc) -- (cpc);
\end{tikzpicture}
\caption{A strict local chain inside the complete lattice of intermediate
logics from Definition~\ref{def:p08-intermediate-lattice}; it is not the whole
superintuitionistic lattice.}
\label{fig:p08-intermediate-lattice}
\end{figure}

\section{Sources and navigation for this chapter}

Generic rules are retrieved from
\hyperref[sheet:lj]{\texttt{sheet:lj}} and
\hyperref[sheet:nj]{\texttt{sheet:nj}}.  The exact inherited IPL adequacy
theorem and its reader-facing sketch are
\hyperref[thm:p08-source-ipl-adequacy]{linked here}; its reciprocal annex item
records the stated omissions.  Separately authored complete proofs retain
their own \texttt{proved=} links.

For Kripke semantics and proof theory, see
\cite{Kripke1965,TroelstraSchwichtenberg2000,TroelstraVanDalen1988}; for
algebraic and topological semantics, see
\cite{RasiowaSikorski1963,McKinseyTarski1946}; for the constructive reading and
formal realizability boundary, see
\cite{Brouwer1908,Heyting1956,Kleene1945}; and for the two intermediate logics,
see \cite{Jankov1968WeakExcludedMiddle,Dummett1959}.
% END HISTORICAL DERIVATIVE 396-681
